% Preamble dùng chung. Biên dịch bằng XeLaTeX (xem latexmkrc).

% --- Ngôn ngữ và phông chữ (Unicode, không cần hack cho tiếng Việt) ---
\usepackage{fontspec}
\usepackage[vietnamese]{babel}
% Chọn font đầu tiên có trên máy: Times New Roman (msttcore-fonts) → Liberation (cùng kích thước,
% có sẵn trên đa số Linux) → TeX Gyre (có trong TeX Live đầy đủ / Overleaf, nạp theo tên file)
% → Latin Modern (mặc định của fontspec, luôn có). Font thực sự được dùng in ra log: tìm "Font chính:".
\newcommand{\chonfont}[4]{% #1 lệnh \set...font, #2 font ưu tiên, #3 font dự phòng, #4 tên file TeX Gyre
  \IfFontExistsTF{#2}{#1{#2}\typeout{Font chính: #2}}{%
    \IfFontExistsTF{#3}{#1{#3}\typeout{Font chính: #3 (thiếu #2)}}{%
      \IfFontExistsTF{#4-regular.otf}{%
        #1{#4}[Extension=.otf, UprightFont=*-regular, BoldFont=*-bold,
               ItalicFont=*-italic, BoldItalicFont=*-bolditalic]%
        \typeout{Font chính: #4 (thiếu #2 và #3)}}{%
        \typeout{Font chính: Latin Modern (thiếu #2, #3, #4)}}}}}
\chonfont{\setmainfont}{Times New Roman}{Liberation Serif}{texgyretermes}
\chonfont{\setsansfont}{Arial}{Liberation Sans}{texgyreheros}
\chonfont{\setmonofont}{Courier New}{Liberation Mono}{texgyrecursor}

% --- Cỡ chữ cố định theo bảng quy ước báo cáo/luận văn ---
% Thân bài 13pt · Mục lớn (section) 14pt · Mục con (subsection) 13pt đậm
% · Chú thích/footnote 10-12pt. Chương và header/listing đặt trực tiếp
% bằng \fontsize tại nơi dùng (xem macros.tex).
\makeatletter
\renewcommand\normalsize{\@setfontsize\normalsize{13}{15.6}} % thân bài: 13pt
\renewcommand\small{\@setfontsize\small{12}{14.4}}           % chú thích bảng/hình: 12pt
\renewcommand\footnotesize{\@setfontsize\footnotesize{10}{12}} % footnote: 10pt
\renewcommand\large{\@setfontsize\large{13}{15.6}}           % tiêu đề mục con (subsection): 13pt
\renewcommand\Large{\@setfontsize\Large{14}{16.8}}           % tiêu đề mục lớn (section): 14pt
\makeatother
\linespread{1.3}  % giãn dòng khoảng 1.5 (chuẩn báo cáo/luận văn)
\setlength{\parindent}{1cm}
\setlength{\parskip}{3pt}

% --- Bố cục: lề trái rộng hơn để đóng gáy ---
\usepackage[a4paper, left=3cm, right=2cm, top=2.5cm, bottom=2.5cm,
            headheight=40pt, headsep=1.0cm]{geometry}
\usepackage{fancyhdr}

% --- Toán, bảng, hình ---
\usepackage{amsmath}
\usepackage{amsthm}     % môi trường định lý / định nghĩa (khai báo ở macros.tex)
\usepackage{unicode-math}
\setmathfont{Latin Modern Math}  % font toán OpenType, chạy mọi cỡ chữ
\usepackage{booktabs, array, tabularx}
\usepackage{graphicx, float, caption}
\usepackage{subcaption}  % nhiều hình con trong một figure
\graphicspath{{figures/}} % \includegraphics tìm ảnh trong figures/
\captionsetup{font=small} % chú thích bảng/hình: 12pt (dùng cỡ \small đã đặt lại ở trên)
\usepackage{xcolor}
\usepackage{colortbl}    % tô màu hàng / ô của bảng
\usepackage{longtable}   % bảng dài tự ngắt trang
\usepackage{siunitx}     % số và đơn vị: \num{}, \SI{}{}, cột S căn theo dấu thập phân
\sisetup{output-decimal-marker={,}, group-separator={.}, group-minimum-digits=4}
\usepackage{algorithm, algorithmic} % khung thuật toán + mã giả
\usepackage{enumerate}
\usepackage[framemethod=tikz]{mdframed}
\usepackage{tikz}
\usetikzlibrary{arrows.meta, positioning}

% --- Mã nguồn ---
\usepackage{listings}

% --- Trích dẫn ---
%% Ghi chú: biblatex khuyến nghị csquotes, nhưng máy chạy lại không cài gói này;
%% thiếu nó chỉ sinh cảnh báo, không ảnh hưởng nội dung.
\usepackage[backend=biber, style=numeric, sorting=none]{biblatex}
\addbibresource{references.bib}
% biblatex không hỗ trợ tiếng Việt: ánh xạ sang tiếng Anh để các chuỗi
% "In:", "et al.", "and" hiện đúng thay vì rơi vào định nghĩa dự phòng hỏng.
\DeclareLanguageMapping{vietnamese}{english}

% --- Siêu liên kết (nạp sau cùng) ---
\usepackage{hyperref}

% Thông tin đề tài: chỉ sửa ở đây, trang bìa và tiêu đề PDF tự cập nhật.
\newcommand{\TenTruong}{Trường Đại học Sài Gòn}
\newcommand{\TenKhoa}{Khoa Công nghệ Thông tin}
\newcommand{\TenMonHoc}{Trí tuệ nhân tạo nâng cao}
\newcommand{\TenDeTai}{Xây dựng và đánh giá Game AI đa cấp độ cho Caro 11×11}
\newcommand{\TenNhom}{Nhóm XX}
\newcommand{\LinkRepo}{https://github.com/SGU-ML25/CaroGame}
\newcommand{\LoaiBaoCao}{Báo cáo đồ án cuối kỳ}
\newcommand{\CommitNop}{chưa chốt}

% Tên file PDF xuất ra cạnh main.tex sau mỗi lần build (không gồm đuôi .pdf; để trống {} thì không xuất).
% Không dùng lệnh LaTeX ở đây, chỉ chữ thường; latexmkrc đọc thẳng dòng này.
\newcommand{\TenFilePDF}{Nhóm XX - Caro 11x11 - TTNTNC}

% Bản nháp / bản nộp. Nháp: dấu mờ "BẢN NHÁP", hiện \goiy và \chualam.
% Nộp (\BanNoptrue): không dấu mờ; còn sót \goiy hoặc \chualam thì build báo lỗi.
\newif\ifBanNop
\BanNopfalse

\newcommand{\GVHD}{(bổ sung)}
\newcommand{\ThoiGian}{TP. Hồ Chí Minh, tháng 12/2026}

% Danh sách sinh viên: mỗi dòng "& Họ tên & MSSV \\" (cột đầu để trống)
\newcommand{\DanhSachSV}{%
  & Huỳnh Thanh Tiến & 3123410374 \\
  & (TV2)            & (MSSV) \\
  & (TV3)            & (MSSV) \\
  & (TV4)            & (MSSV) \\
}

% Macro, màu sắc, thiết lập listings / hyperref / header.

\definecolor{dkgreen}{rgb}{0,0.6,0}
\definecolor{gray}{rgb}{0.5,0.5,0.5}
\definecolor{mauve}{rgb}{0.58,0,0.82}

\hypersetup{
  urlcolor=blue, linkcolor=black, citecolor=black, colorlinks=true,
  pdftitle={\TenDeTai}, pdfauthor={\TenKhoa, \TenTruong},
}

% --- Mã nguồn: listings chỉ dựng bảng ký tự cho mã 0--255. Chữ tiếng Việt có dấu (mã > 255) bị in
% ngay trong khi chữ ASCII còn nằm trong bộ đệm, làm đảo thứ tự ("hoặc" thành "ặhoc"). Hook dưới đây
% đăng ký các chữ đó là chữ cái của listings, giống cách listings tự làm với a--z. ---
\makeatletter
\def\lst@VNletters{258,259,272,273,296,297,360,361,416,417,431,432,7840,7841,7842,7843,7844,7845,7846,7847,7848,7849,7850,7851,7852,7853,7854,7855,7856,7857,7858,7859,7860,7861,7862,7863,7864,7865,7866,7867,7868,7869,7870,7871,7872,7873,7874,7875,7876,7877,7878,7879,7880,7881,7882,7883,7884,7885,7886,7887,7888,7889,7890,7891,7892,7893,7894,7895,7896,7897,7898,7899,7900,7901,7902,7903,7904,7905,7906,7907,7908,7909,7910,7911,7912,7913,7914,7915,7916,7917,7918,7919,7920,7921,7922,7923,7924,7925,7926,7927,7928,7929,9472,9474,9492,9500}
\lst@AddToHook{SelectCharTable}{%
  \@for\lst@vn:=\lst@VNletters\do{%
    \catcode\lst@vn=\active
    \begingroup\lccode`\~=\lst@vn\relax\lccode`\/=\lst@vn\relax
    \lowercase{\endgroup\def~{\lst@ProcessLetter/}}}}
\makeatother

\lstset{
  frame=tb,
  language=Python,
  aboveskip=3mm,
  belowskip=3mm,
  showstringspaces=false,
  columns=flexible,
  basicstyle={\fontsize{10.5}{12.6}\selectfont\ttfamily}, % code/listing: 10.5pt
  numbers=left,
  numberstyle=\tiny\color{gray},
  keywordstyle=\color{blue},
  commentstyle=\color{dkgreen},
  stringstyle=\color{mauve},
  breaklines=true,
  breakatwhitespace=true,
  tabsize=3,
  numbersep=5pt,
  escapeinside={(*@}{@*)},
  captionpos=t
}

% --- Kiểu cho mã LaTeX (dùng trong tài liệu hướng dẫn): \begin{lstlisting}[style=latex] ---
\lstdefinestyle{latex}{language=[LaTeX]TeX, numbers=none, backgroundcolor=\color{gray!8},
  texcsstyle=*\color{blue}, moretexcs={includegraphics,subcaption,caption,label,ref,cite,toprule,midrule,bottomrule,
  multicolumn,rowcolor,cellcolor,num,SI,eqref,goiy,chualam,chuatrienkhai,lstinputlisting,url,href,textcite,
  autoref,cline,addplot,draw,node,STATE,IF,ELSE,ENDIF,FOR,ENDFOR,RETURN,REQUIRE,ENSURE,WHILE,ENDWHILE}}

% --- Định lý / định nghĩa / ví dụ (amsthm), đánh số theo chương ---
\theoremstyle{plain}
\newtheorem{dinhly}{Định lý}[chapter]
\theoremstyle{definition}
\newtheorem{dinhnghia}{Định nghĩa}[chapter]
\newtheorem{vidu}{Ví dụ}[chapter]
\theoremstyle{remark}
\newtheorem*{nhanxet}{Nhận xét}

% --- Tên tiếng Việt cho khung thuật toán ---
\floatname{algorithm}{Thuật toán}
\renewcommand{\algorithmicrequire}{\textbf{Đầu vào:}}
\renewcommand{\algorithmicensure}{\textbf{Đầu ra:}}

% --- Đánh số "Đoạn mã N.M" cho lstlisting, theo chương, giống Bảng/Hình ---
% (listings tự đặt lại số theo chương khi lớp tài liệu có \thechapter)
\renewcommand{\lstlistingname}{Đoạn mã}
\renewcommand{\lstlistlistingname}{Danh sách đoạn mã}

% --- Hộp ghi chú cho phần chưa triển khai ---
\newcommand{\chuatrienkhai}[1]{%
  \begin{mdframed}[backgroundcolor=yellow!15, linecolor=orange, linewidth=1pt]
    \textbf{Chưa triển khai.} #1
  \end{mdframed}}

% --- Ghi chú ngắn cho nội dung sẽ bổ sung ở các tuần sau ---
% --- Gợi ý nội dung cần viết trong sườn bài (xoá khi viết xong mục) ---
% Ở bản nộp (\BanNoptrue trong metadata.tex) cả hai báo lỗi để không nộp sót.
\ifBanNop
  \newcommand{\chualam}[1]{\PackageError{baocao}{Bản nộp còn \string\chualam{#1}}{Viết xong rồi xoá lệnh này.}}
  \newcommand{\goiy}[1]{\PackageError{baocao}{Bản nộp còn \string\goiy{#1}}{Viết xong mục rồi xoá lệnh này.}}
\else
  \newcommand{\chualam}[1]{\textit{(Chưa làm --- #1)}}
  \newcommand{\goiy}[1]{{\color{gray}\itshape [Cần viết: #1]}\par}
\fi

% --- Dấu mờ "BẢN NHÁP" trên mọi trang khi chưa bật \BanNoptrue ---
\ifBanNop\else
  \AddToHook{shipout/background}{%
    \put(0.5\paperwidth,-0.5\paperheight){\makebox(0,0){\rotatebox{45}{%
      \fontsize{90}{90}\selectfont\bfseries\color{gray!12}BẢN NHÁP}}}}
\fi

% --- Logo: bỏ qua êm nếu chưa có file figures/logoITSGU.png ---
\newcommand{\logo}[1]{\IfFileExists{figures/logoITSGU.png}%
  {\includegraphics[#1]{figures/logoITSGU.png}}{}}

% --- Header / footer ---
\setlength{\headheight}{40pt}
\renewcommand{\headrule}{\vskip 2pt\hrule height 0.4pt\vskip 0pt}
\pagestyle{fancy}
\fancyhf{}
\fancyhead[L]{%
  \begin{tabular}{rl}
    \begin{picture}(25,15)(0,0)
      \put(0,-8){\logo{width=8mm,height=8mm}}
    \end{picture} &
    \begin{tabular}[c]{l}
      \normalfont\fontsize{11}{13}\selectfont \TenTruong \\ \TenKhoa % header: 11pt
    \end{tabular}
  \end{tabular}}
\fancyfoot[C]{\thepage}
\renewcommand{\headrulewidth}{0.4pt}
% Trang mở đầu chương dùng cùng kiểu header
\fancypagestyle{plain}{\pagestyle{fancy}}

\setcounter{secnumdepth}{3}
\setcounter{tocdepth}{2}

% --- Tiêu đề chương: căn giữa, in hoa, đậm, một dòng "CHƯƠNG N. TÊN" ---
\makeatletter
\def\@makechapterhead#1{%
  {\parindent\z@ \centering \normalfont\fontsize{17}{20}\selectfont\bfseries % tiêu đề chương: 17pt
   \MakeUppercase{\@chapapp\ \thechapter.\ #1}\par\nobreak
   \vskip 18\p@}}
\def\@makeschapterhead#1{%
  {\parindent\z@ \centering \normalfont\fontsize{17}{20}\selectfont\bfseries % tiêu đề chương: 17pt
   \MakeUppercase{#1}\par\nobreak
   \vskip 18\p@}}
\makeatother

% --- \texttt cho phép xuống dòng sau / _ . để tên file không tràn lề ---
\let\oldtexttt\texttt
\ExplSyntaxOn
\tl_new:N \l__tt_break_tl
\RenewDocumentCommand{\texttt}{m}
  {
    \tl_set:Nn \l__tt_break_tl {#1}
    \tl_replace_all:Nnn \l__tt_break_tl { \_ } { \_ \allowbreak }
    \tl_replace_all:Nnn \l__tt_break_tl { / } { / \allowbreak }
    \tl_replace_all:Nnn \l__tt_break_tl { . } { . \allowbreak }
    \tl_replace_all:Nnn \l__tt_break_tl { - } { - \allowbreak }
    \exp_args:NV \oldtexttt \l__tt_break_tl
  }
\ExplSyntaxOff
\emergencystretch=2em

% --- Từ khoá tiếng Việt cho mã giả (algorithmic) ---
\renewcommand{\algorithmicif}{\textbf{nếu}}
\renewcommand{\algorithmicthen}{\textbf{thì}}
\renewcommand{\algorithmicelse}{\textbf{ngược lại}}
\renewcommand{\algorithmicend}{\textbf{hết}}
\renewcommand{\algorithmicfor}{\textbf{với}}
\renewcommand{\algorithmicdo}{\textbf{làm}}
\renewcommand{\algorithmicwhile}{\textbf{trong khi}}
\renewcommand{\algorithmicreturn}{\textbf{trả về}}

% --- Tên lệnh LaTeX trong bảng / chú thích: \cmd{goiy\{...\}} → \goiy{...} ---
\newcommand{\cmd}[1]{\texttt{\textbackslash #1}}

